\chapter{Trigeminal Neuralgia}
\index{Trigeminal Neuralgia}
\label{sec:Trigeminal Neuralgia}

\section{Overview}
Trigeminal neuralgia (TN, tic douloureux) is a neuropathic pain disorder characterized by brief, electric shock-like, stabbing, or shooting paroxysmal pain in one or more divisions of the trigeminal nerve (typically V2 and V3). Prevalence is 0.01--0.03\%, with a female predominance (2:1) and peak onset after age 50. This neurological disorder affects the central or peripheral nervous system, leading to progressive or episodic symptoms. It significantly impacts quality of life and often requires multidisciplinary care. Neurological disorders affect the central or peripheral nervous system, leading to a wide range of symptoms affecting movement, sensation, cognition, and autonomic function. The nervous system's complexity means that even small disruptions can cause significant functional impairment. Many neurological conditions are chronic and progressive, requiring long-term management and rehabilitation. Advances in neuroimaging and molecular diagnostics have improved understanding and treatment of these conditions.
\section{Causes}
This condition results from disruption of normal trigeminal neuralgia function.

\begin{itemize}[noitemsep]
  \item This condition happens because of neurovascular compression of the trigeminal nerve root entry zone by a blood vessel (most commonly the superior cerebellar artery), leading to focal demyelination and ephaptic transmission
  \item Secondary TN may be caused by multiple sclerosis, tumors, or arteriovenous malformations.
\end{itemize}
\section{Risk Factors}
\begin{itemize}[noitemsep]
  \item Genetic predisposition and family history
  \item Advancing age
  \item Lifestyle factors (diet, exercise, smoking, alcohol)
  \item Environmental exposures
  \item Underlying medical conditions
  \item Medication use
\end{itemize}
\section{Signs and Symptoms}

\subsection{Common symptoms}
\begin{itemize}[noitemsep]
  \item \item You may experience attacks lasting a fraction of a second to 2 minutes triggered by trivial stimuli (light touch, chewing, talking, toothbrushing), with no pain between attacks. Headache and facial pain Weakness or paralysis of muscles Numbness, tingling, or abnormal sensations Vision changes including blurred or double vision Difficulty with balance and coordination Cognitive changes (memory loss, confusion, difficulty concentrating) Speech and language difficulties Seizures or involuntary movements Dizziness and vertigo Fatigue and sleep disturbances

  \item Pain or discomfort in the affected area
  \end{itemize}

\subsection{Emergency warning signs}
\begin{itemize}[noitemsep]
  \item Severe or worsening symptoms of trigeminal neuralgia require immediate medical evaluation
\end{itemize}
\section{Progression and Stages}
The condition may progress through different stages of severity over time. Early stages often involve mild or subtle symptoms that may be overlooked. Moderate stages involve more noticeable symptoms affecting daily function. Advanced stages may involve significant impairment and complications.
\section{Complications}
\begin{itemize}[noitemsep]
  \item Disease progression leading to worsening symptoms
  \item Secondary infections or injuries
  \item Side effects of treatment
  \item Reduced quality of life
  \item Emotional and psychological impact
\end{itemize}
\section{Diagnosis}
Doctors usually diagnose this based on your symptoms and a physical exam with MRI identifying neurovascular compression. Neurological diagnosis combines clinical history and examination findings with neuroimaging (MRI, CT), electrophysiological studies (EEG, EMG, nerve conduction), and laboratory testing of blood and cerebrospinal fluid.

\begin{itemize}[noitemsep]
  \item Comprehensive medical history and physical examination
  \item Laboratory tests to assess organ function
  \item Imaging studies to visualize affected structures
  \item Specialized testing depending on the affected system
  \item Consultation with relevant specialists
\end{itemize}
\section{Prevention}
\begin{itemize}[noitemsep]
  \item Maintain a healthy lifestyle with balanced diet and exercise
  \item Avoid known risk factors
  \item Attend regular medical check-ups for early detection
  \item Follow recommended screening guidelines
\end{itemize}
\section{Treatment Options}
\begin{itemize}[noitemsep]
  \item Treatment options include carbamazepine (200--1200 mg/day) as first-line, with a highly specific response that is itself diagnostic, or oxcarbazepine as an alternative
  \item For refractory cases, surgical options include microvascular decompression (80--90\% long-term pain-free), stereotactic radiosurgery, or through the skin rhizotomy.
  \item Medications to manage symptoms and address underlying mechanisms
  \item Lifestyle modifications and self-care measures
  \item Physical therapy and rehabilitation
  \item Surgical intervention when indicated
  \item Regular monitoring and follow-up care
\end{itemize}
\section{Living with It}
\begin{itemize}[noitemsep]
  \item Follow your treatment plan as prescribed
  \item Maintain regular follow-up with your healthcare provider
  \item Make necessary lifestyle adjustments
  \item Seek support from family, friends, or support groups
  \item Stay informed about your condition
\end{itemize}
\section{Outlook}
\begin{itemize}[noitemsep]
  \item The outlook is generally positive with treatment
  \item Microvascular decompression addresses the underlying cause.
\end{itemize}
